\section{Threats to Validity}
\label{sec:validity}

We state the limits in the body because they bound the central claims.

\smallskip
\noindent\textbf{The effective sample is 19 identities.} The test split's $280$
queries come from $19$ identities visible in two or more photographs across $11$
walls; queries from one wall share material, illumination and one camera pass, so
query count is not sample size. We resample walls and report the
leave-one-wall-out spread in Rank-1. The same resampling discipline applies to
the open-set headline: $\mathrm{DIR}@\mathrm{FAR}{=}0.1$ of $0.768$ is a point
estimate at one operating point on the same sample, so the abstract compares it
to baselines as a one-number structure --- the $0.35$ gap --- and claims
nothing about the precise spread between resampled walls. This benchmark supports
a structural claim about the formulation, but not fine-grained ranking among
baselines. We do not attach a numeric wall-resampled interval to the abstract
point estimates because the score matrices needed to recompute it are not part
of the released artefacts; the effective sample and wall-level bimodality should
therefore be read as a limitation, not as hidden precision.

\smallskip
\noindent\textbf{Scope is controlled, and the control found something.} The
method reads full photographs; the baselines read crops --- dictated by design,
since homography needs wall context. Two controls
(Sec.~\ref{sec:controls}) separate identity cue from field of view, and scope
does not account for the open-set gap ($0.404$ best context against $0.768$).
But the crack-erased control is uncomfortable: inpainting the crack costs only
$0.011$ of that $0.404$, so a context-scope number is partly background-matching,
which we expect to weaken on real revisits. Ranking numbers are additionally
near-duplicate (the adjacent frame answers all $280$ queries,
Sec.~\ref{sec:dataset}); the gap sweep shows the trend is flat, with power
falling to $105$ answerable queries by gap $3$.

\smallskip
\noindent\textbf{The annotation and baseline are circularly coupled.}
Identities were seeded by propagating one click point between consecutive
photographs through an ORB homography (Sec.~\ref{sec:annotation}), and the
geometric baseline's whole mechanism is homography registration. Of the $118$
positive (same-crack) query--gallery pairs on test, $116$ ($98.3\%$) involve an
identity formed by a merge; both-sides-unmerged positives number just $2$.
Thus the clean subgroup needed to test whether geometry helps independently of
annotation provenance does not exist at useful power. The two reported checks
are limited: the procedures are not identical, and merged versus unmerged
positives register at the same observed rate ($100\%$ both ways), but the latter
cannot test whether a biased homography would register differently from an
unbiased one, because both alignments are related procedures. Nor did we run a
sensitivity analysis quantifying how much false-positive mass or registration
selection would be required to manufacture the $0.35$ open-set gap; the saved
score matrices needed for that analysis are not released. Accordingly, the
geometric advantage remains unresolved with respect to circularity, rather than
being demonstrated independent of it. The decisive follow-up is a small,
manually verified identity subset labelled without homography propagation,
reported even if underpowered, together with a bias-sensitivity curve.

\smallskip
\noindent\textbf{Annotations are instrumented, not audited.} The human verdicts
that would turn the provenance of Sec.~\ref{sec:annotation} into a sampled error
rate with a confidence interval have not been collected; the $45\%$ merge rate is
reported, and relative comparison is robust to it since the merge policy hits all
methods identically. All captures come from one site over four half-days, two
handsets and two surface families; smooth plaster produces the categorical SIFT
failure noted in Sec.~\ref{sec:coverage}; cracks are located by a learned
segmenter whose errors propagate into Chamfer agreement unquantified. What would
convert the within-visit benchmark into a cross-visit measurement is a real
second capture session at marked distances and angles; we have not run one.
The synthetic revisit protocol is therefore a controlled check, not a substitute
for multi-visit data, and the paper does not claim cross-visit performance.
